\chapter*{Preface}
\addcontentsline{toc}{chapter}{Preface}

There is a particular moment, familiar to anyone who has really learned a piece
of mathematics, when a formula stops being a string of symbols and becomes a
\emph{picture}. Before that moment you can plug numbers in and get the right
answer; after it, you can close your eyes and \emph{see} why the answer had to
come out that way. This book is written entirely in service of that second kind
of knowing.

Probability is unusually rich in such moments, and unusually generous about
handing them to you, because almost every result in the subject is secretly a
statement about something you can imagine doing: flipping a coin, drawing a card,
waiting for a bus, letting a rumor spread through a town. The mathematics is a
careful bookkeeping of \emph{what tends to happen when chance is in charge}. Our
whole approach is to keep that concrete activity in view at all times -- to let
you watch the dice before we ever write down a distribution.

\section*{What this book is for}

Most probability courses, and many fine textbooks, are organized around results
to be proved and formulas to be applied. You meet the binomial coefficient, then
the binomial distribution, then its mean and variance, and you are asked to
trust that this is all going somewhere. Often it is -- but the trust is doing a
lot of the work, and trust is exactly what evaporates under exam pressure or
three years later when you actually need the idea.

We have written this book on a different bet: that if you genuinely
\emph{understand} where a formula comes from, you will rarely need to memorize
it, and you will be able to rebuild it from scratch when you have forgotten the
details (you will forget the details; everyone does). Understanding is more
robust than memory. It also transfers. The person who truly sees why the law of
large numbers works is the same person who, years later, instantly smells
something wrong in a viral statistic, a clinical-trial summary, or a backtested
trading strategy.

\begin{focusbox}
The single organizing principle of this book: \textbf{intuition first, then
history or mechanism, then a real proof.} We never lead with the theorem. We
lead with the situation that makes the theorem feel inevitable, and only then do
we make it precise. If at any point a derivation feels like it appeared out of
nowhere, that is a bug in our writing, not a gap in your ability.
\end{focusbox}

This is not a ``probability without tears'' book that hides the mathematics.
The proofs are here, and they are real. But a proof can be rigorous without
being dry, and it can be complete without being a wall of symbols. Where a
standard text writes ``it follows that,'' we try to write the sentence that
tells you \emph{why} it follows -- the one line a good teacher says out loud at
the board but that somehow never makes it into print.

\section*{How to read this book}

Read it with a pencil, and read it slowly. Probability punishes skimming more
than almost any other branch of mathematics, because so many of its errors are
\emph{plausible}. A wrong answer in calculus usually looks wrong; a wrong answer
in probability often looks completely reasonable until you simulate it and watch
it fail. (This is why the subject is stuffed with paradoxes, and why we devote
so much space to them.)

A few concrete suggestions:

\begin{itemize}
  \item \textbf{When you meet a new claim, guess first.} Before reading our
    explanation of, say, the birthday problem, write down your gut estimate.
    Being wrong on purpose, and then seeing exactly where your intuition
    misfired, is the fastest way to upgrade that intuition permanently.
  \item \textbf{Treat every box as a different ``voice'' in a conversation.}
    The colored boxes (described below) are not decoration. Each one answers a
    specific question you might be asking at that moment.
  \item \textbf{Do not skip the figures.} They are the argument, not an
    illustration of it. We have drawn each one to make a single idea visible.
  \item \textbf{Simulate.} More on this shortly, but: the surest cure for a
    shaky probabilistic intuition is twenty lines of code and ten thousand
    trials.
  \item \textbf{Do the exercises.} Every chapter ends with a set spanning
    computation, concept, and at least one proof. The conceptual ones are where
    the real learning hides.
\end{itemize}

\begin{marginnotebox}
You do not have to read linearly, but the early chapters genuinely build on one
another. Chapters 1--9 are the load-bearing foundation (sample spaces,
conditioning, random variables, expectation, variance). After that the book
branches: continuous distributions and limit theorems (Chapters
10--22), processes that unfold in time (Chapters 23--29), inference and modern
applications (Chapters 30--35), and a rigorous measure-theoretic capstone
(Chapters 36--38). Two appendices collect the calculus, sums, and linear
algebra you may want to refresh.
\end{marginnotebox}

\section*{The visual language: a guide to the boxes}

Throughout the book, certain kinds of remarks are pulled out into colored
boxes. The point is that \emph{different questions deserve different answers},
and once you learn the color code you can read at two speeds -- following the
main thread, and dipping into exactly the kind of aside you happen to want.
Here is the whole vocabulary, with a tiny example of each.

\begin{intuitionbox}
\textbf{Intuition} (green). The mental picture before the mathematics. Example:
``Variance measures how far a random quantity typically lands from its average
-- it is the spread of the cloud, not its center.''
\end{intuitionbox}

\begin{whybox}
\textbf{Why it works} (blue). The engine room of a proof or formula -- the load-
bearing reason, stated as plainly as we can manage. Example: ``Expectation is
linear because summing and averaging commute; no independence is needed, which
is why $\E[X+Y]=\E X+\E Y$ always holds.''
\end{whybox}

\begin{happeningbox}
\textbf{What is really happening} (deep green). A step back to name the
phenomenon underneath the symbols. Example: ``The central limit theorem is not
about coins or dice in particular; it is about \emph{averaging washing out
detail}, so that almost any source of small independent wiggles produces the
same bell-shaped fog.''
\end{happeningbox}

\begin{historybox}
\textbf{History} (yellow). Where an idea came from, and the human problem that
forced it into existence. Example: probability was born not in a seminar but at
a gambling table, when the Chevalier de M\'er\'e asked Pascal why a betting rule
that had always worked for him suddenly did not.
\end{historybox}

\begin{paradoxbox}
\textbf{Paradox} (red). A result that contradicts the answer almost everyone
gives -- stated sharply, then resolved. Example: switching doors on a game show
doubles your chance of winning, even though, after a door is opened, it
\emph{feels} like a coin flip. We will see exactly why the feeling is wrong.
\end{paradoxbox}

\begin{applicationbox}
\textbf{Application} (magenta). The idea at work in the world -- machine
learning, A/B testing, finance, genetics, signal processing. Example: the same
Bayes' rule that updates a medical diagnosis is the spam filter quietly sorting
your inbox.
\end{applicationbox}

\begin{connectionbox}
\textbf{Connection} (teal). A bridge to the rest of mathematics, often pointing
to the companion \emph{Calculus From the Ground Up} and \emph{Linear Algebra
From the Ground Up}. Example: a probability density is just an integrand, and
expectation is an integral against it -- so everything you know about
integration is suddenly about averages.
\end{connectionbox}

\begin{trapbox}
\textbf{Common trap} (red). A specific, tempting mistake -- placed exactly where
you are most likely to make it. Example: confusing $\P(A \given B)$ with
$\P(B \given A)$, the error at the heart of countless courtroom and medical
blunders.
\end{trapbox}

\begin{walkbox}
\textbf{Walking question} (orange). A question to chew on while you walk away
from the page; we usually answer it shortly after. Example: ``If you keep
flipping a fair coin, must the number of heads and tails eventually equal out?''
(The answer is more surprising than it sounds.)
\end{walkbox}

\begin{labbox}
\textbf{Work it yourself} (purple). A short, hands-on task -- often a
simulation -- where you produce the insight rather than read it.
\end{labbox}

\begin{aestheticbox}
\textbf{Why this matters} (cyan). The longer-range payoff, so that no result
feels like an isolated trophy.
\end{aestheticbox}

There are a few more -- \textsf{Focus}, \textsf{Formal}, \textsf{Reading note},
and \textsf{Pause} -- whose names say what they do. You have already met
\textsf{Focus} and \textsf{Reading note} above. Do not worry about
memorizing the palette; after a chapter or two the colors will mean something to
you without effort, the same way road signs do.

\section*{The arc: from a coin flip to measure theory}

Seeing the whole journey before taking the first step is worth the detour,
because the destination explains the route.

We begin (Chapter 1) with the most elementary object imaginable: a single fair
coin. From it we extract the entire conceptual skeleton of the subject --
outcomes, events, the idea that probability is a kind of \emph{mass} spread over
possibilities. Counting (Chapter 2) lets us handle situations with many equally
likely outcomes, and counting alone takes us surprisingly far. Then comes the
real engine of probabilistic reasoning: \emph{conditioning} (Chapters 4--5).
Learning to update a probability in the light of new information is, more than
any formula, the skill that separates fluent probabilists from the rest.

With those foundations we introduce \emph{random variables} (Chapter 6) --
numbers attached to chance outcomes -- and the two summaries that compress a
whole distribution into a few telling numbers: \emph{expectation} (Chapter 8) and
\emph{variance} (Chapter 9). From there probability goes continuous (Chapters
10--12), where it becomes area under a curve and the bell curve makes its
entrance.

The middle of the book is about \emph{how randomness adds up}. When you combine
many random pieces -- summing them, averaging them, letting time run -- striking
regularities emerge from the chaos. This is the territory of the law of large
numbers and the central limit theorem (Chapters 19--22), arguably the two most
important facts in all of applied mathematics. The figure below is a small
preview: as we average more and more coin flips, the running average is pulled,
as if by gravity, toward the true probability $1/2$.

\begin{figure}[H]
\centering
\begin{tikzpicture}
\begin{axis}[
  width=12cm, height=6.4cm,
  xlabel={number of flips $n$},
  ylabel={running fraction of heads},
  xmin=0, xmax=200, ymin=0, ymax=1,
  ymajorgrids, xmajorgrids, grid style={gray!18},
  legend pos=north east, legend cell align=left,
  title={Order emerging from randomness: a fair coin's running average},
]
\addplot[densityred, thick, dashed, domain=0:200, samples=2] {0.5};
\addplot[sourceblue, thick] coordinates {
(1,1.000)(2,0.500)(3,0.667)(4,0.500)(5,0.600)(6,0.500)(7,0.571)(8,0.500)
(9,0.444)(10,0.500)(12,0.583)(14,0.571)(16,0.563)(18,0.500)(20,0.450)
(24,0.417)(28,0.464)(32,0.531)(36,0.556)(40,0.525)(46,0.543)(52,0.519)
(58,0.483)(64,0.469)(72,0.486)(80,0.512)(90,0.522)(100,0.510)(112,0.491)
(124,0.500)(136,0.507)(150,0.493)(165,0.503)(180,0.506)(200,0.498)
};
\legend{true value $1/2$, running average}
\end{axis}
\end{tikzpicture}
\caption{A single simulated run of 200 fair coin flips. Early on, the running
fraction of heads swings wildly; as $n$ grows it is reeled toward $1/2$. The law
of large numbers (Chapter 20) explains the pull, and the central limit theorem
(Chapter 21) explains exactly how large the remaining wiggles tend to be.}
\label{ch00:lln-preview}
\end{figure}

After the limit theorems we let time enter the picture and study
\emph{processes}: Markov chains, the Poisson process, random walks,
martingales, and Brownian motion (Chapters 23--29). Here probability stops being
a snapshot and becomes a movie. Then we turn the whole apparatus around and use
it for \emph{inference} (Chapters 30--35): given data, what can we legitimately
conclude about the world that produced it? This is where probability meets
statistics, machine learning, and data science.

Finally, in the capstone chapters (36--38), we return to the very beginning --
``what is a probability, really?'' -- and answer it properly, with measure
theory. This is the rigorous foundation that makes everything earlier
unshakable. We have deliberately placed it last. You will appreciate why
Kolmogorov's axioms are shaped the way they are only \emph{after} you have felt,
across thirty-five chapters, the places where naive intuition needs a
guardrail.

\begin{aestheticbox}
Notice the shape of the journey: it is a spiral, not a line. We begin with the
coin flip, travel out through the entire applied subject, and return to the coin
flip to ask what it really was all along. Each pass over a familiar idea is
deeper than the last. If a concept feels only half-understood the first time you
meet it, that is by design -- you will meet it again, equipped to see more.
\end{aestheticbox}

\begin{historybox}
That spiral mirrors the history of the field. Probability spent its first two
centuries (Cardano, Pascal, Fermat, the Bernoullis, Laplace) as a brilliant but
informal art of calculation, applied to games, annuities, and errors of
astronomical measurement. Only in 1933 did Andrey Kolmogorov give it the
measure-theoretic foundation that placed it among the rigorous branches of
mathematics. We follow the same path the subject did: intuition and calculation
first, foundations once we know what they must support. There is no shame in
this order; it is how humanity itself learned probability.
\end{historybox}

\section*{Simulate, then trust, then verify}

If we could press one habit on you, it would be this: \textbf{when in doubt,
simulate.} A computer that flips ten million coins in a second is a remarkably
reliable probability tutor. It has no intuition to fool and no formula
to misremember; it just runs the experiment and reports what happened.

\begin{labbox}
Here, in a few lines of pseudocode, is the entire Figure~\ref{ch00:lln-preview}
experiment: flip a fair coin and report the running fraction of heads after each
flip. We will write code in this spirit throughout the book.
\begin{lstlisting}[language=Python]
import random
heads = 0
for n in range(1, 201):
    heads += (random.random() < 0.5)   # one fair flip
    print(n, heads / n)                # running fraction of heads
\end{lstlisting}
Run it. Watch the printed fraction swing wildly at first, then creep toward
$0.5$ and refuse to leave its neighborhood. (Push the loop to two hundred
thousand flips and it will hug $0.5$ even more tightly.) You have just witnessed
the law of large numbers with your own eyes, before we have proved a single
thing about it.
\end{labbox}

Simulation plays three distinct roles in this book, and keeping them
separate in your mind pays off.

\begin{enumerate}
  \item \textbf{Discovery.} When you do not yet have a guess, simulate to form
    one. The numbers will suggest a pattern; your job is then to explain it.
  \item \textbf{Verification.} When you have derived a formula, simulate to
    check it. If the theory and the simulation disagree, one of them is wrong,
    and finding out which is enormously instructive.
  \item \textbf{Computation.} For problems too hard to solve exactly -- and modern
    applications are full of them -- simulation \emph{is} the answer. Chapter 35
    (Monte Carlo methods) rests entirely on this idea, and it is one of the
    quiet workhorses of modern science.
\end{enumerate}

But notice the order in our section title: simulate, \emph{then trust}, then
verify. There is a temptation, once you have a fast simulator, to stop
thinking and just measure everything. Resist it. A simulation can tell you
\emph{that} something is true to high confidence; only a proof tells you
\emph{why}, and only the why generalizes. A measured probability of $0.4999$
across ten million flips is persuasive, but it is the \emph{argument} that the
true value is exactly $1/2$ -- by symmetry, with no appeal to any experiment at
all -- that you can carry into every other symmetric situation you ever meet.

\begin{happeningbox}
Simulation and proof are not rivals; they are two instruments measuring the same
object. The simulation is your laboratory: it shows you what nature does. The
proof is your theory: it explains why nature had no choice. A mature
probabilist moves fluidly between them -- guessing with the laboratory, securing
with the theory, and trusting neither blindly. Cultivating that two-handed
fluency is, more than any list of distributions, what this book is trying to
teach.
\end{happeningbox}

\begin{paradoxbox}
A warning about ``trust,'' delivered up front so you remember it. Your raw
probabilistic intuition is, frankly, unreliable -- not because you are bad at
it, but because the human brain was simply not built for it. Consider: in a room
of just 23 people, the chance that two share a birthday is already better than
even. Almost everyone guesses the number must be in the hundreds. The intuition
that misleads you here is the same one used by every confident person who has
ever been badly wrong about risk. Our remedy is not ``trust your gut'' -- it is
``train your gut, then verify it.'' By the end of this book your intuition will
be a far better instrument, precisely because you will have spent so much time
catching it in the act of being wrong.
\end{paradoxbox}

Because that birthday claim is the single best advertisement for the whole
enterprise, let us actually do the arithmetic now -- not because you need to
follow every step on a first reading, but so you can see the shape of the kind
of reasoning this book will teach you to perform without flinching.

\begin{example}[The birthday problem, worked once, slowly]
With $n$ people in a room and $365$ equally likely birthdays, it is far easier
to compute the chance that \emph{nobody} shares a birthday and subtract from one.
Line the people up. The first person's birthday is free. The second must dodge
one taken day, so it is ``safe'' with probability $364/365$. The third must dodge
two, giving $363/365$, and so on. Multiplying these step-by-step chances
together (each conditioned on the people before it landing on distinct days),
\[
  \P(\text{all distinct})
  = \frac{365}{365}\cdot\frac{364}{365}\cdots\frac{365-n+1}{365}
  = \prod_{k=0}^{n-1}\Bigl(1-\frac{k}{365}\Bigr).
\]
The probability of \emph{at least one match} is then $1$ minus this product.
Plugging in $n=23$ gives $0.5073$, just over a coin flip. The reason it feels
impossible is that we instinctively count the $23$ people, when what matters is
the number of \emph{pairs} of people, $\binom{23}{2}=253$ -- and $253$ chances
to collide is a great many. Figure~\ref{ch00:birthday} shows the whole curve.
\end{example}

\begin{figure}[H]
\centering
\begin{tikzpicture}
\begin{axis}[
  width=12cm, height=6.2cm,
  xlabel={number of people $n$},
  ylabel={$\P(\text{some shared birthday})$},
  xmin=0, xmax=80, ymin=0, ymax=1.02,
  ymajorgrids, xmajorgrids, grid style={gray!18},
  legend pos=south east, legend cell align=left,
  title={The birthday problem: probability of a shared birthday},
]
\addplot[densityred, thick, dashed, domain=0:80, samples=2] {0.5};
\addplot[sourceblue, very thick] coordinates {
(1,0.000)(5,0.027)(10,0.117)(15,0.253)(20,0.411)(22,0.476)(23,0.507)
(25,0.569)(30,0.706)(35,0.814)(40,0.891)(45,0.941)(50,0.970)(55,0.986)
(60,0.994)(65,0.998)(70,0.999)(80,1.000)
};
\addplot[targetgreen, only marks, mark=*, mark size=2pt] coordinates {(23,0.507)};
\node[anchor=west, font=\footnotesize] at (axis cs:24,0.46) {$n=23$};
\legend{even odds, exact probability}
\end{axis}
\end{tikzpicture}
\caption{The probability that at least two of $n$ people share a birthday,
computed exactly from the formula in the worked example. The curve crosses the
even-odds line (dashed) at the famous $n=23$, marked in green, and is essentially
certain by $n=60$. The steepness early on is the surprise: the count of
\emph{pairs} grows like $n^2$, so the matches pile up far faster than the linear
headcount suggests.}
\label{ch00:birthday}
\end{figure}

\begin{whybox}
Notice what just happened, because it is the rhythm of the entire book. We met a
claim that offends intuition; we did not argue about it, we \emph{computed} it;
and the computation then handed back a \emph{new and better} intuition (``count
the pairs, not the people''). Intuition, calculation, and a sharper intuition --
that loop, run thousands of times, is what turns a beginner into someone fluent.
\end{whybox}

\section*{A note on rigor}

We promise two things about the mathematics. We will not lie to you for
simplicity's sake: when a result holds only under conditions, we say so, and when
a ``proof'' is really a plausibility argument, we flag it clearly and point to
where the genuine proof lives (often in Chapters 36--38). Nor will we hide behind
rigor. A definition that is technically precise but motivationally empty has
failed at its real job, which is to capture an idea you already half-see.

\begin{connectionbox}
This book leans on calculus (for densities, expectations as integrals, and
generating functions) and on linear algebra (for covariance matrices, Markov
chains, and Gaussian vectors). We have written it to be self-contained -- the
appendices gather what you need -- but where a deeper treatment helps, we point
to the companion \emph{Calculus From the Ground Up} and \emph{Linear Algebra
From the Ground Up}, which share this book's philosophy and its box-colored
visual language. The three are meant to be read as a family.
\end{connectionbox}

\section*{Acknowledgments and an invitation}

Every probability book stands on the shoulders of a few great ones, and the
debts of this one are easy to name. From Blitzstein and Hwang we learned that a
rigorous probability text can also be genuinely warm and story-driven. From Ross
we inherited the breadth of worked examples and applications. From Feller, the
conviction that the deep results deserve deep treatment and that the subject's
beauty is part of its content. From Wasserman, the insistence that all of this
must ultimately help us reason from data. Where we have succeeded, much credit
is theirs; where we have stumbled, the fault is ours alone.

Most of all, this book is an invitation. Probability is the mathematics of the
uncertain world -- which is to say, of the only world there is. Learning it well
does not make the future predictable, but it does something better: it makes the
future \emph{intelligible}. You will come to see chance not as the absence of
order but as a different, subtler kind of order, with its own laws as firm as
gravity. Few intellectual pleasures compare to the moment when a tangle of
uncertainty suddenly resolves into a clean, inevitable answer.

One more preview, because it is the result the whole middle of the book builds
toward. We claimed that averaging ``washes out detail.'' Here is what that means
made visible. Take the humble outcome of summing a handful of dice. A single die
is flat -- every face equally likely. Sum two and the distribution becomes a
tent, peaked at $7$. Sum five and it is already a recognizable bell. Nothing
about the die changed; the \emph{shape} emerged purely from adding independent
copies together. That convergence, dramatized in Figure~\ref{ch00:clt}, is the
central limit theorem in miniature.

\begin{figure}[H]
\centering
\begin{tikzpicture}
\begin{axis}[
  width=12cm, height=6.4cm,
  xlabel={sum of the dice (rescaled to a common center)},
  ylabel={relative frequency},
  xmin=-3.2, xmax=3.2, ymin=0, ymax=0.62,
  ymajorgrids, xmajorgrids, grid style={gray!18},
  legend pos=north east, legend cell align=left,
  title={Adding independent dice: a bell curve assembling itself},
]
% one die: flat-ish, wide
\addplot[quotientorange, thick] coordinates {
(-2.0,0.167)(-1.2,0.167)(-0.4,0.167)(0.4,0.167)(1.2,0.167)(2.0,0.167)
};
% two dice: triangular (11 sums, peak 6/36 at the center)
\addplot[targetgreen, thick] coordinates {
(-2.40,0.028)(-1.92,0.056)(-1.44,0.083)(-0.96,0.111)(-0.48,0.139)(0.00,0.167)
(0.48,0.139)(0.96,0.111)(1.44,0.083)(1.92,0.056)(2.40,0.028)
};
% five dice: nearly normal (sampled gaussian-ish)
\addplot[sourceblue, very thick, smooth] coordinates {
(-3.0,0.012)(-2.5,0.035)(-2.0,0.085)(-1.5,0.175)(-1.0,0.300)(-0.5,0.430)
(0,0.500)(0.5,0.430)(1.0,0.300)(1.5,0.175)(2.0,0.085)(2.5,0.035)(3.0,0.012)
};
% limiting normal curve, dashed
\addplot[densityred, thick, dashed, domain=-3.2:3.2, samples=80]
  {0.5*exp(-x^2/1.4)};
\legend{one die (flat), two dice (tent), five dice, limiting bell}
\end{axis}
\end{tikzpicture}
\caption{The distribution of a sum of fair dice, each curve rescaled to share a
center and width. One die is flat; two dice form a triangular tent peaked in the
middle; five dice already hug the smooth limiting bell (dashed). The astonishing
fact, proved in Chapter~21, is that the same bell appears no matter what we are
averaging -- dice, heights, measurement errors, monthly returns -- as long as the
pieces are independent and none dominates.}
\label{ch00:clt}
\end{figure}

\begin{aestheticbox}
If a single picture could justify learning probability, it might be this one. The
bell curve is not a fact about dice; it is a fact about \emph{addition under
uncertainty}. The reason it governs measurement error in physics, sampling error
in polling, and diffusion in finance alike is that all three are sums of many
small independent pushes. Learn why the bell appears once, and you have
understood it everywhere at once.
\end{aestheticbox}

\section*{A few warm-up questions}

You cannot do the real exercises yet -- we have not taught you anything. But the
spirit of this book is that you should always be \emph{guessing}, so here are
five questions to carry into Chapter~1. Do not look anything up. Write a number
or a sentence for each, seal it in an envelope (metaphorically), and check it
against your answer when the relevant chapter arrives. The gap between your guess
now and the truth later is precisely the intuition this book exists to install.

\begin{exercise}[Guess before you learn]
For each, commit to an answer \emph{before} reading on in the book.
\begin{enumerate}
  \item A fair coin is flipped $10$ times. Which exact sequence is more likely:
    \texttt{HHHHHHHHHH} or \texttt{HTHHTHTTHH}? (Careful.)
  \item In the same $10$ flips, are you more likely to get exactly $5$ heads, or
    some number of heads that is \emph{not} $5$?
  \item A test for a rare disease is ``$99\%$ accurate.'' You test positive. Off
    the top of your head, what is the chance you actually have the disease:
    closer to $99\%$, or closer to $10\%$? (This one ruins careers.)
  \item You wait at a bus stop where buses arrive ``on average every $10$
    minutes,'' at random times. You show up at a random moment. Do you expect to
    wait about $5$ minutes, or longer?
  \item Roll a fair die repeatedly until you see a $6$. On average, how many
    rolls will that take?
\end{enumerate}
\end{exercise}

\begin{walkbox}
Take just question 1 on your next walk. Most people feel the second sequence is
``more likely'' because it looks random. But we asked for one \emph{specific}
sequence, and every specific sequence of $10$ flips has the identical probability
$(1/2)^{10}$. The illusion comes from quietly comparing ``ten heads'' against
``the huge \emph{family} of jumbled-looking sequences.'' Holding that distinction
straight -- a single outcome versus a class of outcomes -- is half of all of
probability. Question 2 is its mirror image, and your answers should disagree.
\end{walkbox}

Turn the page. The first thing we are going to do is flip a coin -- and we are
going to take it absolutely seriously.

\begin{focusbox}
\textbf{The compact version of everything above.} Build intuition first.
Distrust your untrained gut, especially when it feels most certain. Simulate to
discover and to check, but prove to understand and to generalize. Read the
boxes, draw the pictures, do the exercises. And remember that the whole
enormous edifice -- from the central limit theorem to Brownian motion to
Bayesian machine learning -- is built, brick by careful brick, on top of a
single fair coin flip.
\end{focusbox}
